\section{Динамика с трением}
\subsection{Базовые задачи на трением}
%%%%%%%%%%%%%%%%%%%%%%%%%%%%%%%%%%%%%%%%%%%%%%%%%%%%
%%%%%%%%%%%%%%%%%%%%%%%%%%%%%%%%%%%%%%%%%%%%%%%%%%%%


 \z{На тело массы $m$, вначале покоившееся на горизонтальной плоскости, в течение времени $\tau$ действует горизонтальная сила $F$. Какое расстояние $L$ пройдет тело за все время движения? Коэффициент трения тела о плоскость $\mu$.}
%
{$s=\frac{F- \mu mg}{2\mu m^2}F\tau^2$}
%%%%%%%%%%%%%%%%%%%%%%%%%%%%%%%%%%%%%%%%%%%%%%%%%%%%
\z{Тело массы $m = 20$ кг тянут с силой $F = 120$ Н по горизонтальной поверхности. Если эта сила приложена под углом $\alpha_1  = 60\degree$ к горизонту, то тело движется равномерно. С каким ускорением $a$ будет двигаться тело, если ту же силу приложить под углом $\alpha_2  = 30\degree$ к горизонту?}
%
{$a=0.8\,\text{м}/\text{с}^2$}
%%%%%%%%%%%%%%%%%%%%%%%%%%%%%%%%%%%%%%%%%%%%%%%%%%%%	
 \z[{\begin{wrapfigure}[5]{r}{0.3\textwidth}%
		\vspace{-12pt}%
		\begin{flushright}%
 			\includegraphics[width=0.3\textwidth]{5.1.3}%
		\end{flushright}%
	\end{wrapfigure}%
	}]%
	{На горизонтальном шероховатом столе помещены грузы $М$ (внизу) и $m$, связанные нитью, переброшенной через неподвижный блок. Коэффициенты трения грузов друг о друга и нижнего груза о поверхность стола одинаковы и равны $\mu$. С какой горизонтальной силой нужно потянуть нижний груз, чтобы тела пришли в движение? Нити горизонтальны.}
 %
{$F=\mu (M+3m)g$}
%%%%%%%%%%%%%%%%%%%%%%%%%%%%%%%%%%%%%%%%%%%%%%%%%%%% 
 \z[{\begin{wrapfigure}[5]{r}{0.3\textwidth}%
		\vspace{-12pt}%
		\begin{flushright}%
 			\includegraphics[width=0.3\textwidth]{5.1.4}%
		\end{flushright}%
	\end{wrapfigure}%
	}]%
	{Брусок массы $М$ лежит на горизонтальной плоскости. На бруске лежит тело массы $m$ (см. рисунок). Коэффициенты трения между телом и бруском, а также между бруском и плоскостью, одинаковы и равны $\mu$. К бруску приложена сила $F$, действующая в горизонтальном направлении.
 	\lpu При каком значении $F_1$ силы $F$ эта система начнет двигаться? \lpuend
 	\lpu При каком значении $F_2$ силы $F$ тело начнет скользить по бруску?  \lpuend
 	\lpu Сила $F$ такова, что тело скользит по бруску. Через какое время $t$ тело упадет с бруска, если длина бруска равна $L$. Размерами тела пренебречь. \lpuend
	}
%
{\puan $F_1= \mu(m+M)g;$~\puan $F_2=2\mu(m+M)g;$~\puan$t= \sqrt{\frac{2LM}{F-2\mu(m+M)g}}$}
%%%%%%%%%%%%%%%%%%%%%%%%%%%%%%%%%%%%%%%%%%%%%%%%%%%% 
 \z{Мальчик, стоя на льду, пытается сдвинуть ящик за привязанную к нему веревку. Масса ящика $100$ кг, масса мальчика $50$ кг. Коэффициент трения ящика о лед $0.3$, а мальчика $0.4$. Под каким минимальным углом к горизонту мальчик должен тянуть веревку, чтобы сдвинуть ящик?}
%
{$\alpha \approx 27 \degree $}
%%%%%%%%%%%%%%%%%%%%%%%%%%%%%%%%%%%%%%%%%%%%%%%%%%%%